\documentclass[11pt]{article}
\usepackage[T1]{fontenc}
\usepackage{lmodern}
\usepackage[margin=1in]{geometry}
\usepackage{amsmath,amssymb,amsthm,mathtools}
\usepackage{booktabs,longtable}
\usepackage{microtype}
\usepackage[hidelinks]{hyperref}
\usepackage{url}
\hypersetup{pdftitle={Quadratic positivity for finite-field q-hit residues},
  pdfauthor={Xiaojian Zeng},
  pdfkeywords={q-hit numbers, finite fields, rook theory, exact enumeration}}
% URL wrapping uses package defaults.
\newtheorem{theorem}{Theorem}
\newtheorem{lemma}[theorem]{Lemma}
\newtheorem{proposition}[theorem]{Proposition}
\DeclareMathOperator{\per}{per}
\DeclareMathOperator{\rank}{rank}
\DeclareMathOperator{\supp}{supp}
\newcommand{\FF}{\mathbb F}
\newcommand{\ZZ}{\mathbb Z}
\newcommand{\QQ}{\mathbb Q}
\newcommand{\NN}{\mathbb Z_{\geq0}}
\newcommand{\hit}{\mathcal H}
\newcommand{\invh}{h^{*-1}}
\title{Quadratic positivity for finite-field $q$-hit residues}
\author{Xiaojian Zeng\\
  \small PBC School of Finance, Tsinghua University, 100190\\
  \small\texttt{xiaojianzeng@foxmail.com}}
\date{}

\begin{document}
\maketitle
\begin{abstract}
The finite-field $q$-hit numbers of an arbitrary board need not be
polynomials in $q$, but their residues modulo low powers of $q-1$ have
well-defined integer coefficients. We prove that their quadratic residue
coefficients are nonnegative for every square board and every hit index.
The proof uses convolution over balanced row and column restrictions.
After division by the classical hit family in this convolution algebra,
the quadratic coefficient has a kernel supported on boards of size at most
four. A patch decomposition of auxiliary permutations proves this
locality bound. Exact enumeration verifies nonnegativity of all remaining
kernel coefficients, with a separate implementation reconstructing the
orbit data and the full finite certificate. The conclusion concerns the
quadratic residue and makes no assumption of global polynomiality.
\end{abstract}

\section{The quadratic residue problem}

Let $B\subseteq[n]\times[n]$, where $n\geq1$, and identify $B$ with its
$0$--$1$ matrix. For a prime power $q$, let $m_i(B,q)$ count the matrices
over $\FF_q$ of rank $i$ whose support is contained in $B$, and put
$M_i(B,q)=m_i(B,q)/(q-1)^i$. Write
\[
 [r]!_q=\prod_{j=1}^{r}(1+q+\cdots+q^{j-1}),\qquad [0]!_q=1.
\]
The matrix $q$-hit numbers of Lewis and Morales
\cite[Eq.~(1.3)]{LewisMorales2020} are defined by
\begin{equation}\label{eq:hit}
 \hit_n(B;q,t):=\sum_{d=0}^{n}H_d(B,q)t^d
 =q^{\binom n2}\sum_{i=0}^{n}M_i(B,q)[n-i]!_q
       \prod_{j=0}^{i-1}(tq^{-j}-1).
\end{equation}
All empty products equal $1$. The normalized counts $M_i$ are integers;
this follows from the diagonal action recalled below, or from
\cite[Proposition~5.1]{LewisEtAl2011}. The entire weight of $M_i$ in
\eqref{eq:hit} is the integer polynomial
\begin{equation}\label{eq:integralweight}
 W_{n,i}(q,t)=q^{\binom n2-\binom i2}[n-i]!_q
                         \prod_{j=0}^{i-1}(t-q^j).
\end{equation}
In particular, the negative powers in the original definition cause no
congruence difficulty.

Set $x=q-1$. Chen and Selover prove that both $M_i(B,x+1)$ and
$H_d(B,x+1)$ have integer polynomial residues modulo $x^6$
\cite[Corollary~3.13]{ChenSelover2026}. We use only their residues modulo
$x^3$. Thus there are unique integers $c_a(B,d)$, $0\leq a\leq2$, such that
\begin{equation}\label{eq:residue}
 H_d(B,x+1)\equiv c_0(B,d)+c_1(B,d)x+c_2(B,d)x^2\pmod{x^3}
\end{equation}
for every $x=q-1$ with $q$ a prime power. This is an integer congruence on
that unbounded set. It is not a Taylor expansion of an assumed polynomial
counting function. To see uniqueness, divide the difference of two
degree-at-most-two representatives by $x^3$. The quotient is an integer
tending to zero along the set, so the difference vanishes at infinitely
many points.

Define $C_{a,n}(B;t)=\sum_{d=0}^{n}c_a(B,d)t^d$.
For the auxiliary size-zero board take $\hit_0=h_0=1$, so
$C_{1,0}=C_{2,0}=0$.

\begin{theorem}\label{thm:main}
For every $n\geq1$, every board $B\subseteq[n]\times[n]$, and every
$0\leq d\leq n$, the quadratic coefficient in \eqref{eq:residue} satisfies
$c_2(B,d)\geq0$.
\end{theorem}

Chen and Selover establish universal linear positivity in their
Theorem~4.10, formulate positivity of the available higher residue
coefficients in Conjecture~4.1, and discuss the quadratic calculation in
Section~5 of \cite{ChenSelover2026}. Theorem~\ref{thm:main} settles the
quadratic case for the full prime-power domain. Its additional ingredient
is a universal identity reducing the quadratic coefficient to kernels of
sizes two, three and four, proved in Theorem~\ref{thm:locality} below.
Positivity of higher residue coefficients is not established here.
The distinction between residues and global polynomiality is necessary:
Stembridge's Fano-board example exhibits characteristic-dependent matrix
counts \cite[Section~7A]{Stembridge1998}.

\section{Balanced restriction convolution}

Put $z=t-1$. Let $J_n$ be the all-ones matrix and set
\begin{equation}\label{eq:classical}
 h_n(B;t)=\per(J_n+zB)
   =\sum_{\pi\in\mathfrak S_n}t^{|\{j:(j,\pi(j))\in B\}|}
   =\sum_{i=0}^{n}r_i(B)(n-i)!z^i,
\end{equation}
where $r_i(B)$ is the number of $i$-edge matchings in $B$. This is the
classical rook--hit identity, also recalled in
\cite[Eq.~(1.1)]{LewisMorales2020}. Indeed, after choosing the $i$ marked
board entries in the permanent expansion, the remaining permutation has
$(n-i)!$ choices. In particular, $h_n(B;t)\in\NN[t]$.

For families $f,g$ of functions on square boards, define
\begin{equation}\label{eq:convolution}
 (f*g)_n(B)=
 \sum_{\substack{R,C\subseteq[n]\\|R|=|C|}}
 f_{|R|}(B[R,C])\,g_{n-|R|}(B[R^c,C^c]).
\end{equation}
Restrictions retain their induced orders. Associativity follows by
partitioning the rows and columns into three corresponding blocks;
commutativity follows by taking complements. The identity has value $1$
in size zero and $0$ in positive sizes. The convolution inverse of $h$ is
\begin{equation}\label{eq:inverse}
 (\invh)_n(B;t)=(-1)^n h_n(B;t).
\end{equation}
For a fixed full permutation, the splittings with a first block of size
$k$ correspond to choosing $k$ of its edges. Its coefficient in the
convolution of the two sides is therefore multiplied by
$\sum_{k=0}^{n}\binom nk(-1)^{n-k}$.

We say $f$ is \emph{local of order at most $k$} if $f=h*A$, where
$A_n=0$ for every $n>k$. For this intermediate definition we allow values
in $\QQ(z)$. The kernel $A=\invh*f$ is unique and satisfies
\begin{equation}\label{eq:recursion}
 A_n(B)=f_n(B)-\sum_{j=0}^{n-1}
 \sum_{\substack{R,C\subseteq[n]\\|R|=|C|=j}}
 A_j(B[R,C])h_{n-j}(B[R^c,C^c]).
\end{equation}

Let $E$ multiply the size-$n$ part of a family by $n$, and let
$L=z\,d/dz$. Both are derivations of \eqref{eq:convolution}.

\begin{lemma}\label{lem:operators}
If $f$ is local of order at most $k$, then $Ef$ and $Lf$ are local of
order at most $k+1$. Multiplication by a scalar in $\QQ(z)$ does not
increase the locality order.
\end{lemma}
\begin{proof}
Marking one edge of a permutation in \eqref{eq:classical} gives
$Eh=h*U$ and $Lh=h*V$, where the only nonzero parts of $U,V$ are
\[
 U_1([b])=1+zb,\qquad V_1([b])=zb\qquad(b\in\{0,1\}).
\]
Consequently $E(h*A)=h*(U*A+EA)$ and
$L(h*A)=h*(V*A+LA)$. Convolution adds support-size bounds, whereas $EA$
and $LA$ have the same support bound as $A$. Scalar multiplication can
be applied directly to $A$.
\end{proof}

Thus a polynomial of total degree $d$ in the commuting operators $E,L$
raises locality by at most $d$. In applications below, each rational
scalar in $z$ is applied \emph{after} the indicated polynomial operator;
it is never commuted through $L$.

\section{Small support components and the patch decomposition}

For an exact matrix support, form its bipartite graph on its nonisolated
row and column vertices. On a connected component $G$ with $r$ rows and
$c$ columns, the diagonal row and column action has a one-dimensional
scalar stabilizer. Every orbit therefore has size $(q-1)^{r+c-1}$.
If $O_D(G,q)$ denotes the number of these orbits of rank $D$, the
component contributes to the normalized rank count the factor
\begin{equation}\label{eq:defect}
 x^e O_D(G,x+1),\qquad
 e=r+c-D-1\geq\max(r,c)-1.
\end{equation}
Ranks and exponents add over components. This is the diagonal-orbit
method of \cite[Section~5.1]{LewisEtAl2011}, in the component form
of \cite[Eq.~(3.2)]{ChenSelover2026}.

An isolated edge has $D=1$, $e=0$ and orbit count $1$; it is the only
component of defect zero. Every other nonempty connected component has
defect at least one. Hence components with more than three rows or more
than three columns disappear modulo $x^3$. For the remaining components
$O_D(G,x+1)$ is a polynomial: this is the small-board polynomiality
consequence stated in Remark~2.3 and used in the proof of Theorem~3.12
of \cite{ChenSelover2026}, based on Stembridge's analysis
\cite{Stembridge1998}. We use polynomiality only for these small
components.

Write the normalized rook residue as
\[
 M_i(B,x+1)\equiv\sum_{a=0}^{2}R_{a,i}(B)x^a\pmod{x^3},
 \qquad
 F_{a,n}(B;z)=\sum_{i=0}^{n}R_{a,i}(B)(n-i)!z^i.
\]
The coefficients $R_{a,i}$ are defined by the same residue convention as
\eqref{eq:residue}.

\begin{lemma}\label{lem:rooklocality}
For $a=0,1,2$, the family $F_a$ is local of order at most $2a$.
Moreover, $F_0=h$.
\end{lemma}
\begin{proof}
Fix an actual labelled collection $G$ of $k$ components other than
isolated edges, with a rank specified on each component. Let $r,c,D$
be its total row count, column count and rank. If its total defect is
$e$, then
\begin{equation}\label{eq:sizebound}
 e=r+c-D-k,\qquad k\leq e,\qquad
 s:=r+c-D=e+k\leq2a
\end{equation}
whenever this collection contributes to order $a$. Its coefficient is
obtained from $x^e$ times the product of its orbit polynomials, and is
independent of the ambient board. The collection is counted as a support,
not as all isomorphisms from a representative, so no automorphism
multiplicity is introduced.

Fix such a support inside $B$, and let $B_{\mathrm{out}}$ be the
rectangular restriction outside all its rows and columns. The rest of
the support consists of an arbitrary matching there. Apart from the
fixed orbit coefficient, its contribution to $F_{a,n}$ is
\begin{equation}\label{eq:corecontribution}
 z^D\sum_{j=0}^{n-D}r_j(B_{\mathrm{out}})(n-D-j)!z^j.
\end{equation}
The factorial arguments are nonnegative, since
$j\leq\min(n-r,n-c)\leq n-D$.

Choose and delete any $D$ core rows and any $D$ core columns. On the
remaining $n-D$ rows and columns form a matrix $W$ whose entry is
$1+zB_{ij}$ when both endpoints are outside $G$, and $1$ otherwise.
The expression in \eqref{eq:corecontribution} equals $z^D\per(W)$,
by expansion in marked entries outside $G$. The retained $r-D$ core
rows and $c-D$ core columns will be called dummy vertices. The deleted
vertices need not correspond to a matching of $G$.

Given an auxiliary permutation in $\per(W)$, let $u$ count its edges
from a dummy row to a dummy column. Its \emph{patch} consists of all
vertices of $G$, together with every outside vertex adjacent in this
permutation to a dummy vertex. The patch has equally many rows and
columns, namely
\begin{equation}\label{eq:patchsize}
 r+(c-D-u)=c+(r-D-u)=s-u\leq s.
\end{equation}
Every permutation edge touching a dummy vertex has weight $1$.
Every outside vertex added to the patch is already used by such an
edge. Thus the remaining edges form a completely unrestricted weighted
permutation on the complementary rows and columns, whose sum is the
corresponding classical hit polynomial $h$.

Conversely, the core, the deleted vertex sets, the partial permutation
touching the dummy vertices, and a complementary permutation reconstruct
the auxiliary permutation uniquely. There are no unused outside
vertices in the patch. All core and partial-permutation data depend
only on the board restricted to the patch; all other dependence is in
the complementary $h$. Averaging over the
$\binom rD\binom cD$ choices of deleted vertices makes this local rule
invariant under row and column relabelling. Each deletion choice has
the same expansion \eqref{eq:corecontribution}, so averaging preserves
the identity. Rational local weights are allowed here.

Summing these patches gives $h*A$ with $A$ supported in sizes at most
$s\leq2a$. Even if $s>n$, only auxiliary permutations with
$s-u\leq n$ can occur, so the argument still applies. At $a=0$ there
are only matchings of isolated edges, with coefficient $1$; their
transform is $h$ by \eqref{eq:classical}.
\end{proof}

The orbit coefficients in this argument can be signed. For example,
the rank-two orbit polynomial of the full $2\times2$ support is $x-1$.
Lemma~\ref{lem:rooklocality} is a locality statement, not a termwise
positivity argument for the orbit expansion.

\section{The locality bound for the quadratic hit coefficient}

For $0\leq i\leq n$ put $m=n-i$, $N=\binom n2$, and define
\begin{align*}
 \alpha(m)&=\frac{m(m-1)}4,&
 \beta(m)&=\frac{m(m-1)(m-2)(9m+13)}{288},\\
 I(i)&=\frac{i(i-1)}2,&
 J(i)&=\frac{i(i-1)(i+1)}6,\\
 T(i)&=\frac{i(i-1)(i-2)(3i-1)}{24},&&\\[-2pt]
 a(n,i)&=N+\alpha(n-i),&
 b(n,i)&=\frac{N(N-1)}2+N\alpha(n-i)+\beta(n-i).
\end{align*}
Here $a$ has total degree at most two and $b$ has total degree at most
four in $(n,i)$. Direct product expansion gives
\[
 [m]!_{1+x}=m!\bigl(1+\alpha(m)x+\beta(m)x^2\bigr)+O(x^3).
\]
For example, dividing the $j$th factor by $j$ gives the coefficients
$(j-1)/2$ and $(j-1)(j-2)/6$; summing these and the cross products
gives the displayed formula. Also,
\[
 t(1+x)^{-j}-1=z-tjx+\frac{tj(j+1)}2x^2+O(x^3).
\]
Consequently the constant, linear and quadratic coefficients of
$W_{n,i}(1+x,t)$ are, respectively,
\begin{align}
 &m!z^i,\nonumber\\
 &m!\bigl(a(n,i)z^i-tI(i)z^{i-1}\bigr),\label{eq:weights}\\
 &m!\bigl(b(n,i)z^i+t(J(i)-a(n,i)I(i))z^{i-1}
                       +t^2T(i)z^{i-2}\bigr).\nonumber
\end{align}
The sums used here are
$\sum_{j=0}^{i-1}j=I(i)$,
$\sum_{j=0}^{i-1}j(j+1)/2=J(i)$, and
$\sum_{0\leq j<k<i}jk=T(i)$.
Terms with negative powers of $z$ have zero multipliers at the relevant
small values of $i$; alternatively, the computation takes place in
$\QQ(z)$.

Since $Lz^i=iz^i$, multiplication of the rook residues by
\eqref{eq:weights} yields the family identity
\begin{equation}\label{eq:operatoridentity}
\begin{split}
 C_2={}&F_2+\left(a(E,L)-\frac{t}{z}I(L)\right)F_1\\
 &+\left(b(E,L)+\frac{t}{z}\bigl(J(L)-a(E,L)I(L)\bigr)
                    +\frac{t^2}{z^2}T(L)\right)h.
\end{split}
\end{equation}
The polynomial operators in the second term have degree at most two;
those in the last term have degree at most four. Lemmas
\ref{lem:operators} and \ref{lem:rooklocality} bound the locality orders
of these three terms by $4$, $2+2$ and $0+4$.

\begin{theorem}\label{thm:locality}
Let $K_2=\invh*C_2$ for the convolution \eqref{eq:convolution}.
Then $K_{2,n}(B;t)$ is an integer polynomial, vanishes for $n=0,1$ and
for every $n>4$, and satisfies
\begin{equation}\label{eq:universal}
 C_{2,n}(B;t)=\sum_{k=2}^{\min(4,n)}
  \sum_{\substack{R,C\subseteq[n]\\|R|=|C|=k}}
  K_{2,k}(B[R,C];t)\,h_{n-k}(B[R^c,C^c];t).
\end{equation}
\end{theorem}
\begin{proof}
Equation~\eqref{eq:operatoridentity} and the locality bounds prove
$K_{2,n}=0$ for $n>4$. Both $C_2$ and $h$ have integer-polynomial
values, so \eqref{eq:inverse} shows that their kernel does too.
Thus the intermediate use of $\QQ(z)$ creates no problem at $t=1$.
In size zero $C_2=0$, and in size one the hit polynomial is independent
of $q$, so $K_{2,0}=K_{2,1}=0$. The equation $C_2=h*K_2$ is now
exactly \eqref{eq:universal}.
\end{proof}

All residue manipulations above are valid on every prime power:
\eqref{eq:integralweight} is an integer-polynomial weight, so replacing
$M_i$ by its residue preserves the congruence. The vanishing above size
four is therefore an algebraic conclusion, not an extrapolation from
small boards.

\section{The exact finite certificate}

\begin{proposition}\label{prop:finite}
For every board of size $k=2,3,4$, every coefficient of
$K_{2,k}(B;t)$ is nonnegative.
\end{proposition}
\begin{proof}[Computer-assisted proof]
We describe the exact finite calculation, including its reconstruction
bound. Choose a spanning tree of a connected support $G$. Diagonal
scaling makes its tree entries $1$, uniquely up to the common scalar
which acts trivially. The remaining
$\nu=|E(G)|-r-c+1$ entries range freely over $\FF_q^\times$. Hence
\[
 0\leq O_D(G,x+1)\leq x^\nu.
\]
Small-board polynomiality, already invoked after \eqref{eq:defect},
implies that this polynomial has degree at most $\nu$. In the range
$r,c\leq3$ one has $\nu\leq4$. Accordingly $\nu+1$ exact counts at
distinct field sizes determine it by rational interpolation; further
field sizes check the result. This is interpolation with a proved
degree bound, not an inference of polynomiality from samples.

There are $253$ connected labelled supports with $1\leq r,c\leq3$
and no isolated vertices. Independent row and column permutations
reduce them to the $28$ classes in Appendix~\ref{app:orbits};
transposition is not identified. All their rank-orbit polynomials are
listed there. For a size-$n$ exact support $S$, its normalized rank
residue, with $y$ recording rank, is computed as
\begin{equation}\label{eq:supportproduct}
 P_S(x,y)=\prod_{G\text{ component of }S}
       \left(\sum_D y^D x^{r(G)+c(G)-D-1}O_D(G,x+1)\right)
                  \pmod{x^3}.
\end{equation}
For the empty support the product is $1$. Components exceeding the
size bound are discarded. For a board $B$, summing $P_S$ over
$S\subseteq B$ gives $\sum_i M_i(B,x+1)y^i$ modulo $x^3$.
We then apply the integral weights \eqref{eq:integralweight} and
extract $C_{2,n}$, and apply \eqref{eq:recursion} to get $K_{2,n}$.

Every board mask, and every coefficient from $t^0$ through $t^n$, is
included. The resulting ranges are as follows.
\begin{center}
\begin{tabular}{rrrrrr}
\toprule
$n$ & Boards & Coefficients & Minimum & Maximum & Negative\\
\midrule
1 & 2 & 4 & 0 & 0 & 0\\
2 & 16 & 48 & 0 & 1 & 0\\
3 & 512 & 2,048 & 0 & 41 & 0\\
4 & 65,536 & 327,680 & 0 & 162 & 0\\
\bottomrule
\end{tabular}
\end{center}
The supplementary programs and full coefficient tables implement these
formulas using exact integer and rational arithmetic. The reconstruction
and comparison checks described next reproduce every entry.
\end{proof}

\begin{proof}[Proof of Theorem~\ref{thm:main}]
By \eqref{eq:universal}, $C_{2,n}$ is a sum of products of the finite
kernels in Proposition~\ref{prop:finite} and classical hit polynomials.
Both factors have nonnegative coefficients in $t$. Therefore
$C_{2,n}(B;t)\in\NN[t]$, which is the required assertion for every $d$.
\end{proof}

\subsection*{Enumeration and verification details}

The support sum in \eqref{eq:supportproduct} is evaluated on the Boolean
lattice of the $n^2$ cells by a subset zeta transform. If an array
initially contains the exact-support values $P_S$, then for each cell
$e$, and each mask $S$ containing $e$, replace its current value by the
sum of that value and the value at $S\setminus\{e\}$. After all cells
are processed the entry at $B$ is $\sum_{S\subseteq B}P_S$.
This is the standard subset-lattice transform described in
\cite[Section~2.2]{BjorklundEtAl2006}. The balanced row--column
convolution used in the proof is defined separately in
\eqref{eq:convolution}.

A second complete evaluation uses embeddings of the small connected
cores, matching dynamic programming on the remaining board, and
\eqref{eq:weights}. It extracts the kernel by the direct alternating
inverse \eqref{eq:inverse}, rather than by recursive subtraction.
The two evaluations agree on all $329{,}780$ quadratic hit coefficients
and all $329{,}780$ kernel coefficients for the $66{,}066$ boards of
sizes one through four. These implementations share the initial orbit
table, so agreement between them alone does not check that table.

A separate verifier reconstructs all $253$ labelled connected supports
without importing either implementation. It chooses its own spanning
trees, computes ranks by explicit minors, and reconstructs the orbit
polynomials by exact rational Vandermonde elimination. It enumerates
$83{,}919$ normalized assignments over fields of sizes
$2,3,4,5,7,8,11,13$. The extension fields use the irreducible polynomials
$u^2+u+1$ and $u^3+u+1$. The resulting polynomials agree with every
entry of the orbit table. Using the integral product weights and
direct permutation enumeration for $h$, it recomputes all three
residue coefficients and the quadratic kernels. All $989{,}340$
residue entries and $329{,}780$ kernel entries agree with the stored
tables, and the reconstructed kernels are checked for nonnegativity
directly.

As additional checks, this verifier compares $2{,}907$ coefficients in
\eqref{eq:weights} with direct product expansion for
$0\leq i\leq n\leq16$, and enumerates every square matrix of sizes
one through three over $\FF_2,\FF_3,\FF_4$. The latter covers
$282{,}701$ matrices and $6{,}300$ exact support--rank comparisons,
using Gaussian elimination instead of the orbit calculation's minors.
These checks concern exact counts, including at $q=2$, where a
congruence modulo $(q-1)^3$ by itself would be vacuous.

Using exact predicates and a separately implemented finite search is
also a verification practice in Zeng, Liu, Ratnavelu and Ong's preprint
on ternary Borsuk and kissing-code models
\cite{ZengEtAl2026Certificates} (Section~4, Lemma~3 and Proposition~1).
Their two clique implementations provide a methodological comparison
for our integer enumeration and cross-checks. The mathematical input
that extends our finite certificate to every board size is specifically
the locality identity \eqref{eq:universal}.

\subsection*{Reproducibility, scope and assistance}

The accompanying \texttt{supplement.zip} contains the exact tables,
the calculation programs, and the separate verifier, with a README
giving commands and data formats. All computations use Python~3's
standard library. No floating-point fitting or probabilistic search
is part of the certificate. Board masks use bit $in+j$ for the cell
$(i,j)$ in zero-based coordinates. The complete proof consists of the
algebraic reduction and the finite certificate together.

The result proves the universal quadratic inequality with the precise
prime-power residue interpretation \eqref{eq:residue}. It does not
establish the full higher-coefficient conjecture or polynomiality of
$H_d(B,q)$ on arbitrary boards.

OpenAI Codex assisted the proof development, programming, source
comparison and preparation of this manuscript. A separate model
context examined the argument and supplied the additional verifier.
This does not constitute external peer review or formal proof-assistant
certification. The deterministic checks share the stated polynomiality
dependency.

\clearpage
\appendix
\section{The connected orbit table}\label{app:orbits}

For a support with $r$ rows and $c$ columns, the integer mask is
$\sum_{(i,j)\in G}2^{ic+j}$, with zero-based indices. Entries below are
$O_D(G,x+1)$; a zero is also shown when rank $D$ exceeds the dimensions.
The table contains one representative under independent row and column
permutations, with transposes retained as separate dimension classes.
It supplies all component polynomials used in
\eqref{eq:supportproduct}; high-defect ranks are removed by the power
of $x$ there. The field reconstruction and degree bound are given in
the proof of Proposition~\ref{prop:finite}.

\begin{longtable}{ccrlll}
\toprule
$r$ & $c$ & Mask & $O_1$ & $O_2$ & $O_3$\\
\midrule
\endhead
1 & 1 & 1 & $1$ & $0$ & $0$\\
1 & 2 & 3 & $1$ & $0$ & $0$\\
1 & 3 & 7 & $1$ & $0$ & $0$\\
2 & 1 & 3 & $1$ & $0$ & $0$\\
2 & 2 & 7 & $0$ & $1$ & $0$\\
2 & 2 & 15 & $1$ & $x-1$ & $0$\\
2 & 3 & 15 & $0$ & $1$ & $0$\\
2 & 3 & 29 & $0$ & $1$ & $0$\\
2 & 3 & 31 & $0$ & $x$ & $0$\\
2 & 3 & 63 & $1$ & $x^{2}-1$ & $0$\\
3 & 1 & 7 & $1$ & $0$ & $0$\\
3 & 2 & 23 & $0$ & $1$ & $0$\\
3 & 2 & 27 & $0$ & $1$ & $0$\\
3 & 2 & 31 & $0$ & $x$ & $0$\\
3 & 2 & 63 & $1$ & $x^{2}-1$ & $0$\\
3 & 3 & 79 & $0$ & $1$ & $0$\\
3 & 3 & 87 & $0$ & $0$ & $1$\\
3 & 3 & 93 & $0$ & $0$ & $1$\\
3 & 3 & 94 & $0$ & $0$ & $1$\\
3 & 3 & 95 & $0$ & $0$ & $x$\\
3 & 3 & 119 & $0$ & $1$ & $x-1$\\
3 & 3 & 127 & $0$ & $x$ & $x^{2}-x$\\
3 & 3 & 221 & $0$ & $1$ & $x-1$\\
3 & 3 & 223 & $0$ & $x$ & $x^{2}-x$\\
3 & 3 & 238 & $0$ & $1$ & $x-1$\\
3 & 3 & 239 & $0$ & $x-1$ & $x^{2}-x+1$\\
3 & 3 & 255 & $0$ & $x^{2}-x+1$ & $x^{3}-x^{2}+x-1$\\
3 & 3 & 511 & $1$ & $x^{3}+2x-3$ & $x^{4}-x^{3}-2x+2$\\
\bottomrule
\end{longtable}

\clearpage
\bibliographystyle{plain}
\bibliography{references}
\end{document}
